\documentclass{article}
\usepackage[T1]{fontenc}
\usepackage{lmodern}
\usepackage{graphicx}
\usepackage{amsmath,amssymb}
\usepackage[a4paper,margin=2cm]{geometry}
\usepackage[absolute,overlay]{textpos}
\setlength{\TPHorizModule}{1mm}
\setlength{\TPVertModule}{1mm}
\usepackage[indonesian]{babel}
\usepackage{hyperref}
\setlength{\emergencystretch}{2em}
% unit: Halaman 5

\begin{document}

% === Halaman 5 (llm) ===

\section*{Penurunan Rumus Enkripsi dan Dekripsi RSA}

\begin{itemize}
    \item Teori yang digunakan: Teorema Euler \[ a^{\phi(n)} \equiv 1 \pmod{n} \]
    \item Syarat:
    \begin{enumerate}
        \item $a$ harus relatif prima terhadap $n$, yaitu $\text{PBB}(a, n) = 1$ $\quad(\text{PBB} = gcd)$
        \item $\phi(n) = \text{Totient Euler} = \text{fungsi yang menentukan berapa banyak dari bilangan-bilangan } 1, 2, 3, \dots, n \text{ yang relatif prima terhadap } n.$
    \end{enumerate}
    
    Contoh: $\phi(20) = 8$, sebab terdapat 8 buah yang relatif prima dengan 20, yaitu $1, 3, 7, 9, 11, 13, 17, 19$.
\end{itemize}

Jika $n = pq$ adalah bilangan komposit dengan $p$ dan $q$ prima, maka
\[ \phi(n) = \phi(p) \phi(q) = (p - 1)(q - 1). \]

\end{document}
